% Capítulo 3. Seguridad de Red y Configuración de Firewalls (ufw / iptables)
\section{Seguridad de Red y Configuración de Firewalls (ufw / iptables)}\label{cap:03_seguridad_red_firewall}
Este capítulo define qué tráfico puede llegar a cada VPS. La regla general es que cada capa solo acepta conexiones de la capa anterior: Internet solo llega al Frontend, el Frontend es el único que llega al Backend y el Backend es el único que llega a la Base de Datos. El puerto 22 (SSH) queda abierto en las tres VPS para la administración con Termius.

La protección se configura en dos capas que trabajan juntas:
\begin{enumerate}
  \item \textbf{ufw} (\emph{Uncomplicated Firewall}), con política de rechazar todo lo entrante salvo lo permitido explícitamente.
  \item Una \textbf{regla en la tabla raw de iptables}, aplicada por los scripts de la plataforma, que descarta las conexiones a los puertos de los contenedores que no vengan de la IP autorizada.
\end{enumerate}

% Política de firewall con Podman
\subsection{Política de firewall con Podman}\label{sec:3_1_politica_firewall}

\subsubsection{ufw: política por defecto y reglas de entrada (IN) y de reenvío (FWD)}
ufw es la interfaz de firewall incluida en Ubuntu. Se configura con una política por defecto y una lista de excepciones:
\begin{lstlisting}
# ufw default deny incoming      # rechazar todo lo que entra
# ufw default allow outgoing     # permitir todo lo que sale
# ufw default deny routed        # rechazar el trafico reenviado (valor por defecto)
\end{lstlisting}

Con Podman hay que distinguir dos tipos de tráfico, porque ufw los filtra en lugares distintos:
\begin{table}[H]
\centering
\caption{Tipos de reglas de ufw}
\small
\begin{tabular}{p{2.6cm}p{4.2cm}p{7.6cm}}
\toprule
\textbf{Tipo} & \textbf{Comando} & \textbf{A qué tráfico se aplica} \\
\midrule
Entrada (\texttt{ALLOW IN}) & \texttt{ufw allow ...} & Conexiones dirigidas a un proceso del propio servidor, como \texttt{sshd} en el puerto 22. Cadena \texttt{INPUT}. \\
Reenvío (\texttt{ALLOW FWD}) & \texttt{ufw route allow ...} & Conexiones que el servidor reenvía a otra red. Es el caso de los contenedores: cuando Podman publica un puerto (\texttt{-p 5432:5432}), redirige la conexión a la IP interna del contenedor (\emph{DNAT}) y el paquete pasa por la cadena \texttt{FORWARD}, no por \texttt{INPUT}. \\
\bottomrule
\end{tabular}
\end{table}
Por eso una regla de entrada como \texttt{ufw allow 5432} por sí sola no basta para un contenedor: también se necesita la regla de reenvío correspondiente (\texttt{ufw route allow}). Del mismo modo, una regla \texttt{deny} de entrada no bloquea el puerto de un contenedor.

\subsubsection{Capa adicional: regla en la tabla raw de iptables}
Como segunda capa, independiente de ufw, los scripts de la plataforma agregan una regla en la cadena \texttt{PREROUTING} de la tabla \texttt{raw}. Esta tabla se evalúa \textbf{antes} que cualquier otra (antes de la redirección de Podman y antes de las reglas de ufw), por lo que el paquete no autorizado se descarta en cuanto llega al servidor, sin importar cómo estén configuradas las demás reglas.

\begin{lstlisting}
iptables -t raw -I PREROUTING -i eth0 -p tcp --dport 5432 ! -s 156.67.26.73 -j DROP
\end{lstlisting}
\begin{table}[H]
\centering
\caption{Partes de la regla raw}
\small
\begin{tabular}{p{3.6cm}p{10.8cm}}
\toprule
\textbf{Parte} & \textbf{Significado} \\
\midrule
\texttt{-t raw} & Tabla que iptables evalúa primero. \\
\texttt{-I PREROUTING} & Inserta la regla en la cadena por la que pasa todo paquete al llegar. \\
\texttt{-i eth0} & Solo el tráfico que entra por la interfaz pública. El tráfico local y el de la red interna de Podman no se toca. \\
\texttt{-p tcp -{}-dport 5432} & Conexiones TCP a ese puerto (o rango, como \texttt{8082:8094}). \\
\texttt{! -s 156.67.26.73} & ``Cuyo origen \emph{no} sea esta IP''. La IP autorizada queda exenta. \\
\texttt{-j DROP} & Descartar el paquete sin responder; desde fuera el puerto parece no existir. \\
\bottomrule
\end{tabular}
\end{table}

La interfaz pública en las VPS de Contabo es \texttt{eth0}:
\begin{lstlisting}
# ip route show default | awk '{print $5; exit}'
eth0
\end{lstlisting}

\nota{En Ubuntu 24.04, iptables y ufw escriben sus reglas en \emph{nftables}, el filtro de paquetes actual del kernel. Por eso todas las reglas (las de ufw, la regla raw y las que crea Podman) pueden verse juntas con \texttt{nft list ruleset}.}

% Reglas perimetrales por nodo
\subsection{Reglas perimetrales por nodo}\label{sec:3_2_reglas_perimetrales}

\begin{table}[H]
\centering
\caption{Resumen de la protección por VPS}
\small
\begin{tabular}{p{2.3cm}p{2.3cm}p{5.2cm}p{4.6cm}}
\toprule
\textbf{VPS} & \textbf{ufw} & \textbf{Permitido} & \textbf{Regla raw} \\
\midrule
Base de datos & Activo & 22 desde cualquiera; 5432 solo desde 156.67.26.73 & 5432 solo desde 156.67.26.73 \\
Backend & Activo & 22 desde cualquiera; puertos de los microservicios solo desde 169.58.62.111 & 8082--8094 solo desde 169.58.62.111 \\
Frontend & Inactivo & Todo (solo \texttt{router-nginx} debe atender en 80/443) & Ninguna \\
\bottomrule
\end{tabular}
\end{table}

% ------------------------------------------------------------
\subsubsection{VPS Base de Datos: puerto 5432 exclusivamente desde la IP del VPS Backend}
\paragraph{ufw.} En el VPS Base de Datos (169.58.62.99):
\begin{lstlisting}
# ufw default deny incoming
# ufw default allow outgoing
# ufw allow 22/tcp
# ufw allow from 156.67.26.73 to any port 5432 proto tcp
# ufw route allow from 156.67.26.73 to any port 5432 proto tcp
# ufw enable
\end{lstlisting}
\begin{table}[H]
\centering
\caption{Reglas de ufw del VPS Base de Datos}
\small
\begin{tabular}{p{5.6cm}p{2cm}p{6.8cm}}
\toprule
\textbf{Regla} & \textbf{Tipo} & \textbf{Motivo} \\
\midrule
\texttt{22/tcp} desde cualquiera & IN & Administración por SSH con Termius. Se agrega \textbf{antes} de \texttt{ufw enable} para no perder la conexión. \\
\texttt{5432/tcp} desde 156.67.26.73 & IN & Permite al Backend llegar al puerto publicado de PostgreSQL. \\
\texttt{5432/tcp} desde 156.67.26.73 & FWD & Permite que esa conexión se reenvíe al contenedor \texttt{defensoria-db}. Sin esta regla el Backend no puede conectarse. \\
Todo lo demás & -- & Rechazado por la política \texttt{deny incoming}. \\
\bottomrule
\end{tabular}
\end{table}

\paragraph{Regla raw.} La aplica el script de la base de datos (capítulo \ref{cap:05_despliegue_base_datos}):
\begin{lstlisting}
# cd /apps/aplicaciones/defensoria/database
# bash podman-compose-bd.sh firewall        # aplica la regla
# bash podman-compose-bd.sh firewall-boot   # la reaplica en cada reinicio
\end{lstlisting}
El script descarta en la tabla raw todo lo que llegue a 5432 (y a 5433, puerto de la antigua base histórica) desde una IP distinta del Backend.

% ------------------------------------------------------------
\subsubsection{VPS Backend: puertos de los microservicios exclusivamente desde la IP del VPS Frontend}
\paragraph{ufw.} En el VPS Backend (156.67.26.73):
\begin{lstlisting}
# ufw default deny incoming
# ufw default allow outgoing
# ufw allow 22/tcp
# for p in 8082 8083 8084 8085 8086 8087 8088 8089 8091 8094; do
    ufw allow from 169.58.62.111 to any port $p proto tcp
  done
# ufw route allow in on eth0 out on podman0
# ufw route allow in on podman0 out on eth0
# ufw route allow in on eth0 out on podman1
# ufw route allow in on podman1 out on eth0
# ufw enable
\end{lstlisting}
\begin{table}[H]
\centering
\caption{Reglas de ufw del VPS Backend}
\small
\begin{tabular}{p{5.2cm}p{1.4cm}p{7.8cm}}
\toprule
\textbf{Regla} & \textbf{Tipo} & \textbf{Motivo} \\
\midrule
\texttt{22/tcp} desde cualquiera & IN & Administración por SSH. \\
8082--8089, 8091, 8094 desde 169.58.62.111 & IN & Un puerto por microservicio con ruta pública: Primer Contacto (8082), auth (8083), quejas (8084), notificaciones (8085), catálogo (8086), admin (8087), revisión (8088), chatbot (8089), Subdefensoría (8091) y denunciado (8094). \\
\texttt{eth0} $\rightarrow$ \texttt{podman0} & FWD & Permite que el tráfico que llega por la interfaz pública se reenvíe a la red de los contenedores. \\
\texttt{podman0} $\rightarrow$ \texttt{eth0} & FWD & Permite que los contenedores salgan a Internet y a otras VPS: conexión a la base de datos (169.58.62.99:5432) y envío de correos por SMTP. \\
\texttt{eth0} $\leftrightarrow$ \texttt{podman1} & FWD & Lo mismo para la segunda red de Podman (\texttt{10.89.0.0/24}). \\
\bottomrule
\end{tabular}
\end{table}

\nota{Las reglas de reenvío \texttt{eth0 $\rightarrow$ podman0/podman1} permiten llegar a \textbf{cualquier} puerto de los contenedores desde cualquier origen. Por eso, en el Backend, quien realmente limita los microservicios al VPS Frontend es la regla raw del siguiente párrafo; las reglas de entrada por puerto documentan qué puertos son públicos para el Frontend.}

\paragraph{Regla raw.} La aplica \texttt{firewall-backend.sh} (capítulo \ref{cap:06_despliegue_backend}):
\begin{lstlisting}
# cd /apps/aplicaciones/defensoria/back
# bash firewall-backend.sh          # aplica la regla
# bash firewall-backend.sh boot     # la reaplica en cada reinicio
\end{lstlisting}
Descarta todo lo que llegue por \texttt{eth0} a los puertos 8082--8094 desde una IP distinta del Frontend. Se limita a \texttt{eth0} a propósito: los microservicios se llaman entre sí por la IP pública (por ejemplo \texttt{156.67.26.73:8085}), pero ese tráfico entra por la red de Podman; una regla sin \texttt{-i eth0} lo bloquearía y las peticiones quedarían colgadas.

\begin{table}[H]
\centering
\caption{Puertos del VPS Backend}
\small
\begin{tabular}{p{1.4cm}p{4.2cm}p{3.2cm}p{5.4cm}}
\toprule
\textbf{Puerto} & \textbf{Servicio} & \textbf{Quién puede entrar} & \textbf{Ruta en nginx} \\
\midrule
22 & SSH & Cualquiera & -- \\
8082 & primer-contacto-service & 169.58.62.111 & \texttt{/api/primer-contacto/} \\
8083 & auth-service & 169.58.62.111 & \texttt{/api/auth/} \\
8084 & quejas-service & 169.58.62.111 & \texttt{/api/quejoso/} \\
8085 & notificaciones-service & 169.58.62.111 & \texttt{/api/notificaciones/} \\
8086 & catalogo-service & 169.58.62.111 & \texttt{/api/catalogos/} \\
8087 & admin-service & 169.58.62.111 & \texttt{/api/admin/} \\
8088 & revision-service & 169.58.62.111 & \texttt{/api/revision/} \\
8089 & chatbot-service & 169.58.62.111 & \texttt{/api/chatbot/} \\
8091 & subdefensoria-service & 169.58.62.111 & \texttt{/api/subdefensoria/} \\
8093 & antecedentes-service (IA) & Solo otros microservicios & Sin ruta pública \\
8094 & denunciado-service & 169.58.62.111 & \texttt{/api/denunciado/} \\
\bottomrule
\end{tabular}
\end{table}

% ------------------------------------------------------------
\subsubsection{VPS Frontend: tráfico público 80 (HTTP) y 443 (HTTPS)}
En el VPS Frontend (169.58.62.111) ufw está \textbf{inactivo} y no hay regla raw:
\begin{lstlisting}
# ufw status
Status: inactive
\end{lstlisting}
Es la única VPS que debe atender a Internet: \texttt{router-nginx} recibe a los usuarios en los puertos 80 y 443, y los cinco frontends (8090--8095) solo deben consumirse a través de él por \texttt{127.0.0.1}.

\paragraph{Configuración recomendada.} Para dejar abiertos únicamente SSH, HTTP y HTTPS, y cerrar a Internet los puertos de los frontends:
\begin{lstlisting}
# ufw default deny incoming
# ufw default allow outgoing
# ufw allow 22/tcp
# ufw allow 80/tcp
# ufw allow 443/tcp
# ufw enable
# IFACE=$(ip route show default | awk '{print $5; exit}')
# iptables -t raw -I PREROUTING -i $IFACE -p tcp --dport 8090:8095 -j DROP
\end{lstlisting}
\texttt{router-nginx} usa la red del host y llega a los frontends por la interfaz local \texttt{lo}, así que ninguna de estas reglas le afecta. La regla raw no sobrevive a un reinicio; para hacerla permanente se agrega a un servicio de systemd como los que crean los scripts del Backend y la Base de Datos.

% ------------------------------------------------------------
\subsubsection{Rutas internas bloqueadas en nginx (403)}
Además del firewall, \texttt{router-nginx} niega desde Internet los endpoints que solo deben usarse entre microservicios. Estas líneas forman parte de \texttt{defensoria.conf} (sección \ref{sec:7_4_defensoria_conf}):
\begin{lstlisting}
location = /api/notificaciones/enviar    { return 403; }
location = /api/notificaciones/registrar { return 403; }
location /api/primer-contacto/ingesta/   { return 403; }
location /api/subdefensoria/ingesta/     { return 403; }
location /api/subdefensoria/acuerdos-conclusion/interno/ { return 403; }
\end{lstlisting}
Los microservicios siguen pudiendo usarlos porque se llaman directamente por \texttt{156.67.26.73:<puerto>}, sin pasar por nginx.

% ------------------------------------------------------------
\subsubsection{Persistencia de las reglas tras reiniciar}
\begin{itemize}
  \item \textbf{ufw}: al ejecutar \texttt{ufw enable} queda habilitado el servicio \texttt{ufw.service}, que carga las reglas guardadas en cada arranque. No requiere nada adicional.
  \item \textbf{Regla raw}: iptables no la guarda por sí mismo. Los comandos \texttt{firewall-boot} (BD) y \texttt{boot} (Backend) de los scripts crean un servicio de systemd que la vuelve a aplicar al encender la VPS:
\end{itemize}
\begin{table}[H]
\centering
\caption{Servicios de systemd del firewall}
\footnotesize
\begin{tabular}{p{2.2cm}p{5.6cm}p{6.4cm}}
\toprule
\textbf{VPS} & \textbf{Servicio} & \textbf{Lo crea} \\
\midrule
Base de datos & \texttt{ufw.service}, \texttt{defensoria-bd-firewall.service} & \texttt{ufw enable}, \texttt{podman-compose-bd.sh firewall-boot} \\
Backend & \texttt{ufw.service}, \texttt{defensoria-backend-firewall.service} & \texttt{ufw enable}, \texttt{firewall-backend.sh boot} \\
\bottomrule
\end{tabular}
\end{table}
\begin{lstlisting}
# systemctl is-enabled ufw defensoria-bd-firewall        # VPS Base de Datos
# systemctl is-enabled ufw defensoria-backend-firewall   # VPS Backend
\end{lstlisting}

\paragraph{Verificación de las reglas de cada VPS.} En cada VPS se listan las reglas de ufw numeradas y la regla raw:
\begin{lstlisting}
# ufw status numbered
# iptables -t raw -L PREROUTING -n -v --line-numbers
\end{lstlisting}
\begin{itemize}
  \item \texttt{ufw status numbered}: lista cada regla con su número, el puerto (\texttt{To}), la acción (\texttt{ALLOW IN} o \texttt{ALLOW FWD}) y el origen (\texttt{From}). El número sirve para borrar una regla con \texttt{ufw delete <número>}.
  \item \texttt{iptables -t raw -L PREROUTING -n -v}: muestra la regla raw y sus contadores \texttt{pkts}/\texttt{bytes}, que indican cuántos paquetes no autorizados ha descartado.
\end{itemize}
\captura{red-01-firewall-bd}{Reglas de ufw y regla raw del VPS Base de Datos}
\captura{red-02-firewall-backend}{Reglas de ufw y regla raw del VPS Backend}
\captura{red-03-firewall-frontend}{Estado del firewall del VPS Frontend}

% Diagnóstico de conectividad
\subsection{Diagnóstico de conectividad}\label{sec:3_3_diagnostico_conectividad}
Estas pruebas sirven para confirmar que el firewall funciona y para diagnosticar cuando un servicio no responde. Mientras no se hayan levantado los contenedores (capítulos 5 a 7), los puertos de la plataforma aparecerán cerrados en todos los casos.

\subsubsection{Verificación de puertos y sockets (ss -tulpn, ufw status, iptables -t raw)}
\texttt{ss} muestra qué puertos están escuchando en el servidor y qué proceso los atiende:
\begin{lstlisting}
# ss -tulpn
\end{lstlisting}
\begin{table}[H]
\centering
\caption{Puertos que deben aparecer en escucha con la plataforma instalada}
\small
\begin{tabular}{lp{10.5cm}}
\toprule
\textbf{VPS} & \textbf{Puertos en escucha (\texttt{LISTEN})} \\
\midrule
Frontend & 22 (sshd), 80 y 443 (nginx de router-nginx), 8090--8095 (frontends) \\
Backend & 22 (sshd), 8082--8089, 8091, 8093 y 8094 (microservicios y modelo de IA) \\
Base de datos & 22 (sshd), 5432 (defensoria-db) \\
\bottomrule
\end{tabular}
\end{table}
Que un puerto aparezca en \texttt{LISTEN} solo indica que el servicio está levantado; quién puede llegar a él lo deciden ufw y la regla raw (sección \ref{sec:3_2_reglas_perimetrales}).

ufw registra las conexiones que rechaza con el prefijo \texttt{[UFW BLOCK]}, útil para saber si el firewall está bloqueando algo que debería pasar:
\begin{lstlisting}
# journalctl -k --since "1 hour ago" | grep "UFW BLOCK" | tail
\end{lstlisting}

\subsubsection{Pruebas de conexión entre nodos y desde Internet}
Con \texttt{nc} (incluido en Ubuntu) se prueba si un puerto remoto acepta la conexión; \texttt{-z} solo comprueba sin enviar datos, \texttt{-v} muestra el resultado y \texttt{-w 3} espera como máximo 3 segundos:
\begin{lstlisting}
# nc -zv -w 3 169.58.62.99 5432      # desde el Backend: debe conectar
# nc -zv -w 3 156.67.26.73 8084      # desde el Frontend: debe conectar
\end{lstlisting}
Desde el equipo de administración (fuera de las VPS), los mismos puertos deben fallar por tiempo de espera:
\begin{lstlisting}
$ nc -zv -w 3 169.58.62.99 5432      # debe fallar (timed out)
$ nc -zv -w 3 156.67.26.73 8084      # debe fallar (timed out)
$ nc -zv -w 3 169.58.62.111 443      # debe conectar
\end{lstlisting}

\begin{table}[H]
\centering
\caption{Matriz de pruebas de conectividad (con la plataforma levantada)}
\small
\begin{tabular}{lllc}
\toprule
\textbf{Desde} & \textbf{Hacia} & \textbf{Puerto} & \textbf{Resultado esperado} \\
\midrule
Internet & Frontend 169.58.62.111 & 80 / 443 & Conecta \\
Internet & Backend 156.67.26.73 & 8084 & Falla \\
Internet & BD 169.58.62.99 & 5432 & Falla \\
Frontend & Backend 156.67.26.73 & 8084 & Conecta \\
Frontend & BD 169.58.62.99 & 5432 & Falla \\
Backend & BD 169.58.62.99 & 5432 & Conecta \\
\bottomrule
\end{tabular}
\end{table}
Un resultado \texttt{succeeded} donde se espera \texttt{timed out} indica que falta una regla en esa VPS; uno \texttt{Connection refused} indica que el firewall deja pasar el tráfico pero el contenedor no está levantado.
\captura{red-04-pruebas-nc}{Pruebas de conexión entre nodos y desde Internet}

